% =======================================================
% 8. TRACCIAMENTO E STATO DEI REQUISITI
% =======================================================
\section{Tracciamento e Stato dei Requisiti}
\label{sec:tracciamento}
Al fine di garantire che l'architettura progettata e le tecnologie scelte soddisfino le funzionalità concordate, si fornisce di seguito la tracciabilità tra le componenti architetturali descritte in questo documento e i requisiti individuati in \link{https://synergo-unit-unipd.github.io/Second-Brain/docs/PB/doc_tecnici/doc_esterni/analisi_dei_requisiti/analisi_dei_requisiti.pdf}{Analisi dei Requisiti} (Sezione 4 -- Requisiti), a cui si rimanda per la descrizione testuale completa di ciascun requisito.

\subsection{Tabella Tracciamento Componenti-Requisiti}

\renewcommand{\arraystretch}{1.3}
\begin{xltabular}{\textwidth}{|>{\raggedright\arraybackslash\small}p{3.4cm}|X|>{\raggedright\arraybackslash\small}p{2.6cm}|}
	\caption{Tracciamento Componenti $\rightarrow$ Requisiti} \label{tab:tracciamento-componenti} \\
	\hline
	\rowcolor{primarycolor!10}
	\textbf{Componente} & \textbf{Responsabilità} & \textbf{Requisiti coperti} \\
	\hline
	\endfirsthead

	\multicolumn{3}{c}{{\bfseries \tablename\ \thetable{} -- Continua dalla pagina precedente}} \\
	\hline
	\rowcolor{primarycolor!10}
	\textbf{Componente} & \textbf{Responsabilità} & \textbf{Requisiti coperti} \\
	\hline
	\endhead

	\hline
	\multicolumn{3}{|r|}{{Continua nella pagina successiva...}} \\
	\hline
	\endfoot

	\hline
	\endlastfoot

	\texttt{MarkdownContent-}\newline\texttt{Editor},\newline \texttt{EditorView}/\newline\texttt{EditorController} & Scrittura libera, copia/taglia/incolla, formattazione (grassetto, corsivo, sottolineato, barrato, citazione, intestazione), link, elenchi e tabelle & R1--R42 \\ \hline
	\texttt{EditCommand} (famiglia),\newline \texttt{CommandHistory} & Annullamento e ripristino delle modifiche & R43-F-O, R44-F-O \\ \hline
	\texttt{EditorView}\newline (modalità di visualizzazione) & Visualizzazione esclusiva editor/render\textsubscript{G} o affiancata (Split View\textsubscript{G}) & R45--R47 \\ \hline
	\texttt{AIRouter}, \texttt{AIService},\newline \texttt{AIOperationFactory},\newline operazioni Strategy & Riassunto, traduzione, riscrittura, Distant Writing e analisi Sei Cappelli tramite LLM & R48, R50, R52, R53, R55, R57, R59, R61, R63, R65, R67 \\ \hline
	\texttt{ErrorHandlers},\newline gerarchia \texttt{AIDomainError} & Notifica degli errori di comunicazione con i servizi AI & R49, R51, R54, R56, R58, R60, R62, R64, R66, R68 \\ \hline
	\texttt{AIRequestModel},\newline \texttt{AIPanelView},\newline \texttt{AIController} & Accettazione/rifiuto/rigenerazione della proposta, interruzione manuale, stato di avanzamento continuo & R69--R74 \\ \hline
	\texttt{NoteModel},\newline \texttt{NoteService}/\newline\texttt{NoteServiceProxy} & Salvataggio e caricamento della nota in formato Markdown dal file system locale & R75-F-O, R76-F-O \\ \hline
	\texttt{ExportRouter},\newline esportatori per formato (Template Method) & Esportazione della nota in PDF, HTML e JSON & R77-F-O \\ \hline
	\texttt{App.vue} & Richiesta di conferma prima di chiudere/navigare fuori dall'applicazione con modifiche non salvate; apertura dei link dell'anteprima in una nuova scheda o segnalazione di URL non valido & R87-F-O, R88-F-O \\ \hline
	Suite di test\newline automatici, pipeline CI & Copertura di test ($\geq$80\%\ sui moduli principali), test di integrazione con i servizi LLM, gestione non bloccante degli errori & R1-Q-O, R2-Q-O, R3-Q-O, R10-Q-O, R11-Q-O \\ \hline
	python-dotenv,\newline \texttt{OpenAIAdapter} & Gestione delle credenziali di accesso al servizio LLM esclusivamente lato Backend, mai esposte al client o nei log & R8-Q-O \\ \hline
	Architettura\newline Client-Server,\newline adapter LLM & Esecuzione via browser, standard API OpenAI, singolo modello LLM per richiesta, repository pubblico & R1-V-O, R2-V-O, R3-V-O, R5-V-O, R6-V-O, R7-V-O \\ \hline
	Scelta libera del\newline framework Frontend & Vue.js adottato senza vincolo imposto dal capitolato & R11-V-O \\ \hline
	Docker Compose,\newline Dockerfile & Consegna come archivio distribuibile con configurazione Docker, senza deploy su infrastruttura server esterna & R12-V-O \\ \hline
	Rendering reattivo\newline Vue.js, stato\newline \texttt{AIRequestState} & Aggiornamento anteprima entro 1s (anche per note estese), tempo di avvio dell'applicazione, indicatore di caricamento continuo & R1-P-O, R2-P-O, R3-P-O, R4-P-O \\ \hline
\end{xltabular}

I requisiti opzionali R4-V-Op e R78--R86 (persistenza remota, gestione utenti) non compaiono in questa tabella poiché il gruppo ha deciso di non implementarli in questa fase (si veda Sezione~7.2 per lo stato dettagliato e la motivazione): non esiste quindi alcun componente architetturale a cui tracciarli.

\begin{minipage}{\textwidth}
    \small 
    \vspace{0.2cm} 
    \textit{Nota:} R8-V-O, R9-V-O e R10-V-O (obbligo di produrre rispettivamente Specifica Tecnica, Piano di Qualifica e Manuale Utente) sono requisiti di natura documentale, soddisfatti dall'esistenza dei documenti stessi: non è quindi individuabile un componente architetturale a cui tracciarli, e non compaiono nella tabella precedente.
\end{minipage}

\subsection{Stato di Soddisfacimento dei Requisiti}
La tabella seguente riassume, per macro-gruppo, lo stato di implementazione dei requisiti funzionali, di qualità, di vincolo e prestazionali rispetto all'architettura descritta in questo documento. I gruppi ricalcano la struttura di \link{https://synergo-unit-unipd.github.io/Second-Brain/docs/PB/doc_tecnici/doc_esterni/analisi_dei_requisiti/analisi_dei_requisiti.pdf}{Analisi dei Requisiti}, Sezione 4; il numero di requisiti obbligatori (100), desiderabili (5) e opzionali (15) è riepilogato nello stesso documento, Sezione 6 -- Riepilogo.

I requisiti opzionali relativi alla persistenza remota e alla gestione utenti (R4-V-Op, R7-Q-Op, R8-Q-Op, R78--R86) sono stati valutati e volutamente esclusi dallo scope\textsubscript{G} della release corrente: il gruppo ha scelto di concentrare le risorse disponibili sul consolidamento e sulla qualità dei requisiti obbligatori. Sono pertanto riportati con stato \textit{Non implementato}.

\renewcommand{\arraystretch}{1.3}
\begin{xltabular}{\textwidth}{|>{\centering\arraybackslash}p{2.6cm}|X|c|c|}
	\caption{Stato dei Requisiti} \label{tab:stato-req} \\
	\hline
	\rowcolor{primarycolor!10}
	\textbf{Codici} & \textbf{Descrizione} & \textbf{Tipo} & \textbf{Stato} \\
	\hline
	\endfirsthead
	
	\multicolumn{4}{c}%
	{{\bfseries \tablename\ \thetable{} -- Continua dalla pagina precedente}} \\
	\hline
	\rowcolor{primarycolor!10}
	\textbf{Codici} & \textbf{Descrizione} & \textbf{Tipo} & \textbf{Stato} \\
	\hline
	\endhead

	\hline
	\multicolumn{4}{|r|}{{Continua nella pagina successiva...}} \\
	\hline
	\endfoot

	\hline
	\endlastfoot

	R1--R4 & Scrittura, copia, taglia, incolla di testo nell'editor. & Obbligatorio & Soddisfatto \\ \hline
	R5--R28 & Attivazione / applicazione / disattivazione / rimozione delle formattazioni grassetto, corsivo, sottolineato, barrato, citazione e intestazione. & Obbligatorio & Soddisfatto \\ \hline
	R29--R31 & Inserimento, modifica e rimozione di link ipertestuali. & Obbligatorio & Soddisfatto \\ \hline
	R32--R35 & Inserimento, modifica, disattivazione e rimozione di elenchi puntati e numerati. & Obbligatorio & Soddisfatto \\ \hline
	R36--R42 & Inserimento e modifica di tabelle Markdown, con validazione dei parametri. & Obbligatorio & Soddisfatto \\ \hline
	R43, R44 & Annullamento (Undo) e ripristino (Redo) dell'ultima modifica. & Obbligatorio & Soddisfatto \\ \hline
	R45--R47 & Visualizzazione esclusiva editor, render o affiancata (Split View). & Obbligatorio & Soddisfatto \\ \hline
	R48, R49 & Generazione riassunto tramite LLM e notifica di errore. & Obbligatorio & Soddisfatto \\ \hline
	R50, R51 & Traduzione (Inglese, Francese, Spagnolo) tramite LLM e notifica di errore. & Obbligatorio & Soddisfatto \\ \hline
	R52 & Traduzione in Tedesco come estensione facoltativa. & Desiderabile & Soddisfatto \\ \hline
	R53, R54 & Riscrittura del testo tramite LLM e notifica di errore. & Obbligatorio & Soddisfatto \\ \hline
	R55, R56 & Distant Writing da prompt utente e notifica di errore. & Obbligatorio & Soddisfatto \\ \hline
	R57--R68 & Analisi critica secondo le sei prospettive dei Cappelli di De Bono e relative notifiche di errore. & Obbligatorio & Soddisfatto \\ \hline
	R69--R71 & Accettazione, rifiuto e rigenerazione della proposta AI. & Obbligatorio & Soddisfatto \\ \hline
	R72, R73 & Visualizzazione continua dello stato di avanzamento e interruzione manuale della richiesta AI. & Obbligatorio & Soddisfatto \\ \hline
	R74 & Visualizzazione della proposta AI in un modale dedicato. & Obbligatorio & Soddisfatto \\ \hline
	R75, R76 & Salvataggio e caricamento della nota in formato Markdown su file system locale. & Obbligatorio & Soddisfatto \\ \hline
	R77 & Esportazione della nota in PDF, HTML e JSON. & Obbligatorio & Soddisfatto \\ \hline
	R87, R88 & Conferma prima di uscire con modifiche non salvate; apertura dei link dell'anteprima (nuova scheda o segnalazione URL non valido). & Obbligatorio & Soddisfatto \\ \hline
	R78--R86 & Creazione/eliminazione nota, registrazione/autenticazione utente, gestione note remote. & Opzionale & Non implementato \\ \hline
	R1-Q-O, R2-Q-O, R3-Q-O & Copertura di test ($\geq$80\% sui moduli principali), versionamento Git, test di integrazione con i servizi LLM. & Obbligatorio & Soddisfatto \\ \hline
	R8-Q-O, R10-Q-O, R11-Q-O & Gestione delle credenziali LLM esclusivamente lato Backend; continuità dell'editing locale e gestione non bloccante degli errori in caso di indisponibilità dei servizi esterni. & Obbligatorio & Soddisfatto \\ \hline
	R4-Q-Op, R5-Q-Op, R9-Q-Op & Test di integrazione, autenticazione e hashing delle password del modulo opzionale di persistenza remota. & Opzionale & Non implementato \\ \hline
	R6-Q-D & Adozione di una licenza open source permissiva. & Desiderabile & Soddisfatto \\ \hline
	R7-Q-D & Criteri minimi di usabilità dell'editor (contrasto, leggibilità, focus da tastiera). & Desiderabile & Soddisfatto \\ \hline
	R12-Q-D & Architettura modulare che consenta la sostituzione del provider LLM senza impatti significativi. & Desiderabile & Soddisfatto \\ \hline
	R13-Q-D & Assenza di dipendenza da funzionalità non standard dei browser. & Desiderabile & Soddisfatto \\ \hline
	R1-V-O -- R3-V-O, R5-V-O -- R7-V-O & Esecuzione via browser standard, interazione con i modelli LLM tramite API compatibile OpenAI, singolo modello per richiesta, repository pubblico. & Obbligatorio & Soddisfatto \\ \hline
	R8-V-O -- R10-V-O & Produzione di Specifica Tecnica, Piano di Qualifica e Manuale Utente. & Obbligatorio & Soddisfatto \\ \hline
	R11-V-O & Libertà di scelta del framework Frontend (Vue.js). & Obbligatorio & Soddisfatto \\ \hline
	R12-V-O & Consegna come archivio Docker distribuibile, senza deploy su infrastruttura server esterna. & Obbligatorio & Soddisfatto \\ \hline
	R4-V-Op & Persistenza remota tramite un database server-side. & Opzionale & Non implementato \\ \hline
	R13-V-Op & Predisposizione per il deploy su server esterno, in aggiunta all'archivio Docker. & Opzionale & Non implementato \\ \hline
	R1-P-O & Aggiornamento dell'anteprima Markdown entro 1 secondo dalla modifica. & Obbligatorio & Soddisfatto \\ \hline
	R2-P-O & Visualizzazione continua e reattiva dello stato di caricamento durante l'elaborazione AI. & Obbligatorio & Soddisfatto \\ \hline
	R3-P-O & Tempo di avvio dell'applicazione entro 3 secondi in condizioni di rete standard. & Obbligatorio & Soddisfatto \\ \hline
	R4-P-O & Rendering dell'anteprima senza degrado percettibile anche per note estese ($>$10.000 caratteri). & Obbligatorio & Soddisfatto \\ \hline
	R5-P-O & Timeout di sicurezza per le richieste AI (soglia elevata, non percepibile in condizioni normali); unico controllo effettivo sui tempi affidato all'interruzione manuale (R73-F-O). & Obbligatorio & Soddisfatto \\ \hline
	R6-P-Op & Caricamento della lista delle note remote entro 3 secondi. & Opzionale & Non implementato \\ \hline
\end{xltabular}
